\documentclass[11pt]{article}
\usepackage{amsmath,amssymb,amsthm}
\usepackage[margin=1in]{geometry}
\usepackage{booktabs}
\usepackage[colorlinks=true,urlcolor=blue]{hyperref}

\title{A localized $K$-Chevalley / $K$-Pieri formula for $LG(n,2n)$:\\
what I hold, at what grade, and one trap}
\author{Clio}
\date{2026-10-22}

\newcommand{\LG}{\mathrm{LG}}
\newcommand{\OG}{\mathrm{OG}}
\newcommand{\GP}{\mathsf{GP}}
\newcommand{\GQ}{\mathsf{GQ}}

\begin{document}
\maketitle

\begin{center}
\fbox{\begin{minipage}{0.92\textwidth}
\textbf{Author:} Clio \quad\textbf{Date:} 2026-10-22 \quad\textbf{Recipient:} Rick
(\texttt{grandparick20@gmail.com})\\
\textbf{Commit:} \texttt{a1b5214} (\texttt{clio-vega/proofs})\\
\textbf{In reply to:} your UID 832, \S6 --- \emph{``do you know a localized
$K$-Chevalley/$K$-Pieri formula for $\LG(n,2n)$?''}
\end{minipage}}
\end{center}

\section{Short answer}

Yes for the Pieri half, with a named owner; partially for the Chevalley half, and
there the best pointer I have is one I hold only at relay grade. But the thing I most
want you to take from this note is not a citation --- it is \textbf{a scope trap that
sits exactly on your side of the problem} (\S3).

This is a report on \emph{my reading state}, not a mathematical result. Every item
below is graded. A grade here records \textbf{provenance depth}, not correctness and
not coverage.

\section{What I hold}

\begin{center}
\begin{tabular}{@{}llp{6.2cm}@{}}
\toprule
\textbf{Item} & \textbf{My grade} & \textbf{What it gives you} \\
\midrule
Ikeda--Naruse, \texttt{arXiv:1112.5223},
\emph{$K$-theoretic analogues of factorial Schur $P$- and $Q$-functions}
& \textbf{verified-quote} & \S8 is \emph{equivariant $K$-theory of the maximal
isotropic Grassmannians}; \S8.3 (\texttt{ssec:K\_T}) is $K_T(G_n)$ and its Schubert
basis --- i.e.\ the \emph{localized} picture, which is where your
$\Psi_{\kappa,0} = \GQ$-at-fixed-points result already lives. \\
\addlinespace
Buch--Ravikumar, \emph{Pieri rules for the $K$-theory of cominuscule Grassmannians}
& \textbf{named, ID not held} & The $K$-Pieri rule. $\LG(n,2n)$ is cominuscule in
type $C_n$, so it is in scope. See \S4 on why I am not giving you an arXiv number. \\
\addlinespace
Arroyo, \texttt{arXiv:2511.05734}
& agent-summary & A re-derivation of the Buch--Ravikumar Pieri rule via different
tableaux. \emph{Relayed, not verified by me.} \\
\addlinespace
Mihalcea--Naruse--Su, \texttt{arXiv:2312.17200},
\emph{Chevalley formulae for the motivic Chern classes of Schubert cells and for the
stable envelopes}
& agent-summary & The closest thing I have to a \emph{localized Chevalley stated in
the stable-envelope language} we both work in. \emph{Relayed, not verified by me} ---
in particular I have \textbf{not} checked that its type-$C$ / isotropic case is
actually written out. \\
\addlinespace
Halacheva--Knutson--Zinn-Justin, \texttt{arXiv:1811.07581},
\emph{Restricting Schubert classes to symplectic Grassmannians using self-dual puzzles}
& \textbf{deep-read}, on disk & Cohomological, not $K$-theoretic --- but it is
restriction-to-fixed-points done by puzzles, so if you want the localized statement in
a combinatorial model this is the template. \\
\addlinespace
Chiu--Marberg, \texttt{arXiv:2111.08993} & \textbf{verified-quote} & States which
$\GP/\GQ$ structure-constant problems are open. \\
\bottomrule
\end{tabular}
\end{center}

\section{The trap --- and it is on your side}

$\GQ \leftrightarrow \LG(n,2n)$ (Lagrangian, type $C$); $\GP \leftrightarrow
\OG(n,2n+1)$ (maximal orthogonal, type $B$). These are \emph{different spaces}, and
the literature is \textbf{not symmetric between them}.

Read first-hand from Marberg \texttt{arXiv:2211.01092} (LaTeX source on my disk),
L2999--3013, Prop.\ \texttt{bialgebras-thm}, verbatim:

\begin{quote}
\emph{``First, the products $\GP_\lambda \GP_\mu$ and $\GQ_\lambda \GQ_\mu$ always
expand as finite linear combinations of $\GP$- and $\GQ$-functions. For the
$\GP$-functions, this was first shown in [CTY]; for other proofs, see [HKPWZZ \S4],
[Marberg2021 \S1.2], or [PechenikYong \S8]. For the $\GQ$-functions, the desired
finiteness property is [LM2022, Thm 1.6].''}
\end{quote}

So: \textbf{four independent proofs on the $\GP$ (orthogonal) side; exactly one on the
$\GQ$ (Lagrangian) side.} Your side rests on a single theorem --- Lewis--Marberg,
\emph{Combinatorial formulas for shifted dual stable Grothendieck polynomials},
Thm 1.6.

The specific way this bites: \textbf{Clifford--Thomas--Yong is the wrong side.} CTY's
combinatorial $K$-theoretic jeu-de-taquin rule is for $\OG(n,2n+1)$. It is the paper
everyone reaches for, it is genuinely the canonical shifted $K$-LR rule, and it does
\textbf{not} cover $\LG(n,2n)$. If you pick it up expecting your space you will get a
rule for the other one.

(One caveat I am recording rather than suppressing: I hold a \emph{measured
bibliography disagreement} on CTY's scope. Ikeda--Naruse's 2013 \texttt{bibitem} reads
$\OG(2,2n+1)$; Marberg's biblatex reads $\OG(n,2n+1)$, Crelle 2014. Both witnesses are
\emph{other people's} bibliographies, neither is the paper. I have not resolved it, so
do not let $\OG(n,2n+1)$ harden on my say-so.)

\section{Why there is no arXiv number against Buch--Ravikumar}

Because I do not hold one. The name occurs about a hundred times across my reading
corpus and my sources index carries no identifier for the paper itself --- only
Arroyo's re-derivation of it.

I am flagging this rather than supplying a plausible number because of something that
happened to me yesterday: a \emph{wrong} arXiv ID for Motegi survived review and
reached GitHub, and it survived precisely because every neighbouring field in the
citation was correct. A citation slot is graded by its neighbours, and its neighbours
cannot grade it. The venue I want to write --- Crelle, 2012 --- is \textbf{recall, not
a resolved lookup}, so treat it as a search hint and not a reference.

\section{What is open, in case that is what you actually need}

If your use is a Pieri/Chevalley \emph{expansion}, the above is enough. If your use is
\emph{positivity of structure constants} in this family, the frontier is sharper and
closer than the Pieri literature suggests. From Marberg \texttt{2211.01092}
L3118--3139 (read in full):

\begin{itemize}
\item \textbf{Known:} $\GP_\lambda \GP_\mu \in m\GP$ and $\GQ_\lambda \GQ_\mu \in
m\GQ$, equivalently $a^\nu_{\lambda\mu}, b^\nu_{\lambda\mu} \ge 0$
[LM2022, Thm 1.6]. Combinatorial rules for $a^\nu_{\lambda\mu}$: [CTY, Thm 1.2] or
[Pechenik--Yong, \S8].
\item \textbf{Open} (Conj.\ 4.35/4.36): $\hat a^\nu_{\lambda\mu} \ge 0$ and
$\hat b^\nu_{\lambda\mu} \ge 0$, equivalently $\Delta(\GP_\lambda) \in m\GP \otimes
m\GP$ and $\Delta(\GQ_\lambda) \in m\GQ \otimes m\GQ$. These are the
\textbf{coproduct} constants.
\item \textbf{Open} (Problem 4.37): combinatorial interpretations of the remaining
integers, \emph{``outside some special cases considered in [BuchRavikumar, LM2022]''}.
\end{itemize}

Note the last line: \textbf{the Buch--Ravikumar Pieri rule is itself one of the named
``special cases'' where the answer is known.} That tells you how far the Pieri rule
reaches and where it stops.

And one sentence of Marberg's that is directly about the object you asked for ---
L3139, verbatim:

\begin{quote}
\emph{``We do not know how to leverage the geometric interpretation of $\GP_\lambda$
and $\GQ_\lambda$ in [IkedaNaruse, \S8.3] to prove these properties.''}
\end{quote}

That is the single-author hub of this literature declaring the route
\emph{equivariant $K$-theory Schubert classes on maximal isotropic Grassmannians
$\Rightarrow$ coproduct positivity} to be open, by name. If your $K0$ result really is
$\Psi_{\kappa,0} = $ Ikeda--Naruse $\GQ$ at fixed points with no shift, then you are
standing in \S8.3. I think that is worth knowing before you spend a slot on the Pieri
rule.

\section{What I did not do}

I did not open Buch--Ravikumar, Mihalcea--Naruse--Su, or Arroyo. I did not verify that
\texttt{2312.17200} treats type $C$. I did not resolve the CTY scope disagreement. The
two items I read first-hand are \texttt{1112.5223} (v3 LaTeX e-print, 2026-10-21) and
\texttt{2211.01092} (LaTeX on disk); everything else above is index metadata and is
labelled as such.

\end{document}
